\documentclass[11pt,a4paper]{book}

% --- Encoding & Schrift ---
\usepackage[utf8]{inputenc}
\usepackage[T1]{fontenc}
\usepackage{lmodern}
\usepackage{microtype} % Bessere Typografie und Randausgleich

% --- Seitenlayout ---
\usepackage[margin=2.5cm]{geometry}

% --- Mathepakete ---
\usepackage{amsmath,amssymb,amsthm}

% --- Sprache ---
\usepackage[ngerman]{babel}

% --- Links im PDF ---
\usepackage{hyperref}

\usepackage{csquotes}
\usepackage[backend=biber, style=numeric]{biblatex}
\addbibresource{literatur.bib}


% --- Umgebungen für Aufgaben & Lösungen ---
\newcounter{aufgabe}[chapter]
\renewcommand{\theaufgabe}{\thechapter.\arabic{aufgabe}}

\newenvironment{aufgabe}[1][]{%
  \refstepcounter{aufgabe}%
  \par\noindent\textbf{Aufgabe \theaufgabe}%
  \ifx#1\empty\else\textbf{ (#1)}\fi%
  \par\medskip
}{\medskip}

\newenvironment{loesung}{%
  \par\noindent\textbf{Lösung zu Aufgabe \theaufgabe}\par\medskip
}{\medskip}

\begin{document}

\title{Lösungsbuch zu \\ \textit{Analysis 1} \\ von Otto Forster und Florian Lindemann}
\author{Nicolas Matheisen}
\date{\today}
\maketitle

\tableofcontents
\setcounter{tocdepth}{0}

\chapter{Vollständige Induktion}

\begin{aufgabe}[{\cite[Seite~20, Aufgabe~1.1]{Forster2023}}]
Man beweise die Summenformel für alle $n \in \mathbb{N}$:
\[
  \sum_{k=1}^{n} k^2 = \frac{n(n+1)(2n+1)}{6}.
\]
\end{aufgabe}

\begin{loesung}
Induktions-Anfang: $n=1$

\[
  \sum_{k=1}^{1} k^2 = 1^2 = 1 \cdot 1 = 1
\]

\[
  \frac{1(1+1)(2\cdot1+1)}{6} = \frac{1\cdot2\cdot3}{6} = \frac{6}{6} = 1
\]

Induktions-Voraussetzung:

\[
  \exists n \in \mathbb{N}: \sum_{k=1}^{n} k^2 = \frac{n(n+1)(2n+1)}{6}
\]

Induktions-Schritt: $n \to n+1$

zu zeigen:

\begin{align*}
    \sum_{k=1}^{n+1} k^2 
    &= \frac{(n+1)(n+1+1)(2\cdot(n+1)+1)}{6} \\
    &= \frac{(n+1)(n+2)(2n+3)}{6}
\end{align*}

\begin{align*}
    \sum_{k=1}^{n+1} k^2 
    &= \left(\sum_{k=1}^{n} k^2\right) + (n+1)^2 \\
    &\stackrel{\text{IV}}{=} \frac{n(n+1)(2n+1)}{6} + (n+1)^2 \\
    &= \frac{n(n+1)(2n+1)}{6} + \frac{(n+1)^2}{1} \\
    &= \frac{n(n+1)(2n+1)}{6} + \frac{6(n+1)^2}{6} \\
    &= \frac{n(n+1)(2n+1)+6(n+1)^2}{6} \\
    &= \frac{(n+1) \cdot [n(2n+1) + 6(n+1)]}{6} \\
    &= \frac{(n+1) \cdot [2n^2+1n+6n+6]}{6} \\
    &= \frac{(n+1) \cdot [2n^2+7n+6]}{6}
\end{align*}

$2n^2+7n+6$ in die Grundform bringen.

\[
2n^2+7n+6 = 2\left(n^2+\frac{7n}{2}+3\right)
\]

\begin{align*}
    p &= \frac{7}{2}, \quad q = 3 \\
    x_{1,2} &= -\frac{p}{2} \pm \sqrt{ \left(\frac{p}{2}\right)^2 - q } \\
    x_{1,2} &= -\frac{7}{4} \pm \sqrt{ \left(\frac{7}{4}\right)^2 - 3 } \\
    x_{1,2} &= -\frac{7}{4} \pm \sqrt{ \frac{49}{16} - \frac{48}{16} } \\
    x_{1,2} &= -\frac{7}{4} \pm \frac{1}{4} \\
    x_1 &= -\frac{6}{4} = \frac{3}{2} = -1,5 \\
    x_2 &= -2
\end{align*}

\begin{align*}
  2n^2 + 7n + 6 &= 2 \cdot (n - x_1) \cdot (n - x_2) \\
  &= 2 \cdot (n + \frac{3}{2})(n + 2) \\
  &= (2n + 3)(n + 2)
\end{align*}

\[
  \frac{(n+1) \cdot [2n^2+7n+6]}{6} = \frac{(n+1)(n+2)(2n+3)}{6}
\]

\end{loesung}

\printbibliography
\end{document}